%-------------------------
% Career-Ops LaTeX CV Template — "clean" (A4, ATS-safe vector fonts)
%------------------------

%%%%%%%%%%%%%%%%%%%%%%%%%%%%  Preamble %%%%%%%%%%%%%%%%%%%%%%%%%%%%
\documentclass[a4paper,11pt]{article}
\usepackage{latexsym}
\usepackage[empty]{fullpage}
\usepackage{titlesec}
\usepackage{marvosym}
\usepackage[usenames,dvipsnames]{color}
\usepackage{verbatim}
\usepackage{enumitem}
\usepackage[hidelinks]{hyperref}
\usepackage{fancyhdr}
\usepackage[english]{babel}
\usepackage{tabularx}
\usepackage[T1]{fontenc}
\usepackage{lmodern}
\input{glyphtounicode}
\pdfgentounicode=1

% fullpage's default margin is not 0.5in; geometry sets the exact figure
% without touching anything above. Appended, not substituted — the packages
% and their order above are unchanged.
\usepackage[a4paper,margin=0.5in]{geometry}

% \pdfgentounicode=1 above is a pdfTeX-only primitive: it does nothing under
% tectonic (XeTeX-based), and microtype's \DisableLigatures — the usual
% engine-agnostic fix — is ALSO pdfTeX-only and hard-errors under XeTeX. The
% actual fix for ligature-safe text extraction under every engine lives in
% lib/latex-escape.mjs's escapeLatex(): it breaks every f+{f,i,l} pair with
% an empty TeX group before the text ever reaches this template, which is
% engine-independent because it acts at the TeX lexer level, not the font.

\pagestyle{fancy}
\fancyhf{}
\fancyfoot{}
\renewcommand{\headrulewidth}{0pt}
\renewcommand{\footrulewidth}{0pt}
\urlstyle{same}
\raggedbottom
\raggedright
\setlength{\tabcolsep}{0in}

% Sections formatting — ONE command, nothing else controls heading spacing.
\titleformat{\section}{
  \vspace{-4pt}\scshape\raggedright\large
}{}{0em}{}[\color{black}\titlerule \vspace{-5pt}]

%%%%%%%%%%%%%%%%%%%%%%%%%%%%  Commands  %%%%%%%%%%%%%%%%%%%%%%%%%%%%
\newcommand{\resumeItem}[1]{
  \item\small{#1 \vspace{-2pt}}
}

\newcommand{\resumeSubheading}[4]{
  \vspace{-2pt}\item
    \begin{tabular*}{0.97\textwidth}[t]{l@{\extracolsep{\fill}}r}
      \textbf{#1} & #2 \\
      \textit{\small#3} & \textit{\small #4} \\
    \end{tabular*}\vspace{-3pt}
}

% 2-argument title/date row, matching the shared JS renderer's contract
% (build-cv-latex.mjs emits exactly {name}{dates} here for every .tex
% template — base, CJK, or a named pack). A project's tech-stack line and its
% visible link line are appended as free text by that same function right
% after this call, inside the same \item — not passed through the macro's
% arguments, so this macro never needs to accept more than 2.
\newcommand{\resumeProjectHeading}[2]{
  \vspace{-2pt}\item
    \begin{tabular*}{0.97\textwidth}[t]{l@{\extracolsep{\fill}}r}
      \textbf{#1} & \textit{\small #2} \\
    \end{tabular*}\vspace{-2pt}
}

\newcommand{\resumeSubItem}[1]{\resumeItem{#1}\vspace{-4pt}}
\newcommand{\resumeSubHeadingListStart}{\begin{itemize}[leftmargin=0.15in, label={}]}
\newcommand{\resumeSubHeadingListEnd}{\end{itemize}}
\newcommand{\resumeItemListStart}{\begin{itemize}[topsep=4pt, label=\textbullet]}
\newcommand{\resumeItemListEnd}{\end{itemize}\vspace{2pt}}

%%%%%%%%%%%%%%%%%%%%%%%%%%%%  Heading  %%%%%%%%%%%%%%%%%%%%%%%%%%%%
\begin{document}

\begin{center}
    \textbf{\Huge \scshape {{NAME}}} \\ \vspace{2pt}
    {{TITLE_LINE}}\small {{PHONE}} $|$
    \href{{{EMAIL_URL}}}{\underline{{{EMAIL_DISPLAY}}}} $|$
    {{LOCATION}} \\ \vspace{1pt}
    \small \href{{{LINKEDIN_URL}}}{\underline{{{LINKEDIN_DISPLAY}}}} $|$
    \href{{{GITHUB_URL}}}{\underline{{{GITHUB_DISPLAY}}}} $|$
    \href{{{PORTFOLIO_URL}}}{\underline{{{PORTFOLIO_DISPLAY}}}}
\end{center}

% Section order for this template: Summary, Work Experience, Projects,
% Education, Technical Skills. Candidate's explicit choice (2026-09-14),
% reversing the earlier Education-first order: with a current role and two
% substantial projects on the page, the evidence of doing the work outranks the
% credential, and a hiring manager reads the bullets first. Markers carry no
% ordering meaning of their own, so moving these blocks is purely a matter of
% physical order in this file.
%
% The Additional block (merged certifications + awards line) was removed: the
% entries it carried were an unglossed internal university award and a 2024
% intro certificate, both superseded by the rest of the page. buildAdditional()
% still runs and its ADDITIONAL key is simply unused — build-cv-latex.mjs only
% fails on placeholders left UNRESOLVED in the template, never on unused keys.
% Restore by pasting back a \section{Additional} block holding the ADDITIONAL
% placeholder. Written without the braces on purpose: substitution is a plain
% text replace that does not know what a LaTeX comment is, so a real placeholder
% here would have the built block spliced in after the `%` and every line of it
% past the first would escape the comment into live document body.

% The Summary block is a single substitution, not a marker-delimited section:
% buildSummary() emits the \section heading and the text together, or an empty
% string, so a blank summary leaves no bare header. Do not wrap it in a
% `%%%% ... %%%%` banner — that is the shape stripEmptySections() looks for, and
% summary is a scalar it cannot evaluate.
{{SUMMARY}}

%%%% Experience %%%%
\section{Work Experience}
  \resumeSubHeadingListStart
{{EXPERIENCE}}
  \resumeSubHeadingListEnd

%%%% PROJECTS %%%%
% "Projects", not "Personal Projects": the entries here include an
% examiner-reviewed rebuild and a delivery to a live client, and filing those
% under "personal" reads them down to hobby work.
\section{Projects}
\resumeSubHeadingListStart
{{PROJECTS}}
\resumeSubHeadingListEnd

%%%% Education %%%%
\section{Education}
  \resumeSubHeadingListStart
{{EDUCATION}}
  \resumeSubHeadingListEnd

%%%% Technical Skills %%%%
\section{Technical Skills}
% No extra \vspace here, unlike the base template's Skills block: the
% \titleformat above already applies \vspace{-5pt} after every section's
% rule. Stacking a second negative vspace pulled this section's first line
% up into its own rule (verified visually -- the two overlapped).
\begin{itemize}
[leftmargin=0.15in, label={}]\small{\item{
{{SKILLS}}
}}
\end{itemize}

%%%% END %%%%

\end{document}
